\documentclass[11pt]{article}
\makeatletter\def\input@path{{../estilo/}}\makeatother
\usepackage{didistribuidores}
\renewcommand{\documentocorto}{Bitácora · Etapa 4}
\newcommand{\evidencia}[1]{\lstinputlisting[style=salida]{pruebas/evidencia/#1}}

\begin{document}
\portada{Bitácora de pruebas}{Etapa 4: resiliencia}%
  {Circuit breaker, reintentos seguros, huella de la petición, recuperación automática y pruebas de caídas}%
  {Rama \cod{04-resiliencia}}{18 de noviembre de 2026}

\tableofcontents
\newpage

\section{Qué se construyó en esta etapa}

Hasta la etapa 3 el sistema ya tenía tiempos límite y mensajería con reintentos. Faltaba qué pasa
cuando una pieza falla \textbf{seguido}: cada petición seguía llamando al servicio caído y esperando su
tiempo límite, así que una falla de repartidores se volvía una espera lenta y generalizada en el núcleo.
La etapa 4 agrega las defensas que faltaban:

\begin{itemize}
  \item \textbf{Circuit breaker} (\cod{app/resiliencia.py}): tras 3 fallos seguidos, el núcleo deja de
        llamar a repartidores y responde 503 al instante; a los 10 s deja pasar una llamada de prueba.
  \item \textbf{Reintentos seguros}: un solo reintento, solo cuando la petición no llegó a enviarse o el
        servicio respondió 502, 503 o 504.
  \item \textbf{Deduplicación completa}: una \cod{Idempotency-Key} reutilizada con otro pedido se rechaza
        con 422 (huella SHA-256 de la petición).
  \item \textbf{Recuperación automática}: los pools de conexiones comprueban cada conexión antes de usarla.
  \item \textbf{Estado operativo}: \cod{GET /estado} muestra el circuito y los eventos pendientes.
  \item \textbf{Interfaces}: el panel web (\cod{/panel/}) y la aplicación de escritorio
        \cod{Proyecto/interfaz/interfaz.py}.
\end{itemize}

\begin{nota}[Esta bitácora sí es contra Docker real]
Todas las salidas de esta bitácora son de una corrida real del sistema completo (siete contenedores,
Docker Desktop 29.7.2 en Windows 11) el 18 de noviembre de 2026. Cada prueba es un script en
\cod{documentacion/bitacora-4/pruebas/} y su salida se guardó, tal cual, en
\cod{pruebas/evidencia/}. Las pruebas encontraron dos diferencias con lo que decía el reporte técnico
de la etapa; las dos se documentan en la sección~\ref{sec:hallazgos} y una se corrigió en el código.
\end{nota}

\subsection{Requisitos de la rama \cod{04-resiliencia}}
\begin{tabularx}{\linewidth}{@{}l l X@{}}
\toprule
\textsf{\textbf{Requisito}} & \textsf{\textbf{Estado}} & \textsf{\textbf{Dónde}} \\
\midrule
Timeout con reembolso & Etapa 1 & Relojes de pago y asignación (bitácora 1, sección 9) \\
Timeout HTTP entre servicios & Etapa 3 & 1 s para conectar y 3 s para leer \\
Deduplicación & Completada ahora & Huella de la petición (prueba 1) \\
Circuit breaker & Nuevo & Pruebas 2, 3 y 4 \\
Reintentos & Nuevo en la llamada a repartidores & Sección \ref{sec:reintentos} \\
Pruebas de caídas & Matriz completa & Sección \ref{sec:matriz} \\
\bottomrule
\end{tabularx}

\section{Cómo funcionan las defensas}

\subsection{Las capas, de afuera hacia adentro}

\begin{center}
\begin{tikzpicture}[every node/.style={font=\sffamily\small}]
  \node[draw=azulsd, fill=azulsd!6, rounded corners=3pt, minimum width=12.5cm, minimum height=3.9cm] (b) {};
  \node[anchor=north west, text=azulsd, font=\sffamily\small\bfseries] at (b.north west) {\ Circuit breaker: ¿está abierto? sí $\rightarrow$ 503 al instante, sin llamar};
  \node[draw=acento, fill=white, rounded corners=3pt, minimum width=11.3cm, minimum height=2.7cm, below=0.6cm of b.north] (r) {};
  \node[anchor=north west, text=acento, font=\sffamily\small\bfseries] at (r.north west) {\ Reintento: máximo 1, solo si es seguro};
  \node[draw=grisTexto, fill=grisfondo, rounded corners=3pt, minimum width=10cm, minimum height=1.5cm, below=0.55cm of r.north] (t) {};
  \node[anchor=north west, text=grisTexto, font=\sffamily\small\bfseries] at (t.north west) {\ Tiempo límite: 1 s para resolver y conectar, 3 s para leer};
  \node[draw, fill=white, rounded corners=3pt, below=0.6cm of t.north] {Llamada HTTP a \texttt{repartidores}};
\end{tikzpicture}
\end{center}

Cada capa resuelve un problema distinto: el tiempo límite evita esperar para siempre; el reintento
salva los fallos momentáneos; el breaker evita seguir golpeando un servicio que ya se sabe caído. Los
reintentos van \textbf{dentro} del breaker: una llamada con reintento cuenta como un solo éxito o un solo
fallo y no llena la cuenta el doble de rápido.

\subsection{El circuit breaker}

\begin{center}
\begin{tikzpicture}[node distance=3.2cm, every node/.style={font=\sffamily\small},
  est/.style={draw, thick, rounded corners=4pt, minimum width=2.6cm, minimum height=0.9cm},
  fl/.style={-Stealth, thick}]
  \node[est, draw=verdeok, fill=verdeok!8] (c) {CERRADO};
  \node[est, draw=rojo, fill=rojo!8, right=of c] (a) {ABIERTO};
  \node[est, draw=ambar, fill=ambar!10, right=of a] (s) {SEMIABIERTO};
  \draw[fl] (c) -- node[above,font=\sffamily\scriptsize]{3 fallos seguidos} (a);
  \draw[fl] (a) -- node[above,font=\sffamily\scriptsize]{pasan 10 s} (s);
  \draw[fl] (s) to[bend left=35] node[below,font=\sffamily\scriptsize]{la prueba falla} (a);
  \draw[fl] (s.north) to[bend right=25] node[above,font=\sffamily\scriptsize]{la prueba sale bien} (c.north);
\end{tikzpicture}
\end{center}

\begin{tabularx}{\linewidth}{@{}l X X@{}}
\toprule
\textsf{\textbf{Estado}} & \textsf{\textbf{Qué significa}} & \textsf{\textbf{Qué hace con una llamada}} \\
\midrule
CERRADO & Todo normal & La deja pasar; cada fallo suma uno, un éxito reinicia la cuenta \\
ABIERTO & Hubo 3 fallos seguidos & La rechaza al instante, sin llamar \\
SEMIABIERTO & Pasó la pausa de 10 s & Deja pasar \textbf{una sola} llamada de prueba; las demás se rechazan mientras tanto \\
\bottomrule
\end{tabularx}

\textbf{Qué cuenta como fallo.} El tiempo agotado, la conexión rechazada, el nombre que no se resuelve,
un 5xx y una respuesta ilegible. Un \textbf{404 cuenta como éxito}: el servicio contestó bien, y que el
repartidor no exista es una respuesta. Si contara como fallo, alguien que consulta repartidores
inexistentes podría abrir el circuito y dejar sin servicio a todos.

\textbf{Detalles de la implementación.} El breaker es seguro entre hilos (los cambios de estado van
dentro de un \cod{Lock}), registra cada transición en el log y recibe el reloj como parámetro, lo que
permitió probarlo con un reloj falso. Solo vive en \cod{app/clientes/repartidores.py}: la lógica de
pedidos solo ve que a veces el cliente lanza \cod{ServicioNoDisponible}, igual que en la etapa 3.

\subsection{Los reintentos}\label{sec:reintentos}

\begin{tabularx}{\linewidth}{@{}X l X@{}}
\toprule
\textsf{\textbf{Situación}} & \textsf{\textbf{¿Reintenta?}} & \textsf{\textbf{Por qué}} \\
\midrule
Conexión rechazada o tiempo de conexión agotado & Sí & La petición nunca se envió; con dos copias, el reintento puede caer en la sana \\
Respuesta 502, 503 o 504 & Sí & Son errores de ``no pude atenderte ahora'' \\
El nombre no se resuelve & No & No hay ninguna copia viva; reintentar solo duplica la espera (ver prueba 4) \\
Tiempo de lectura agotado & No & Ya se esperó el máximo; reintentar duplicaría la espera \\
500 u otro error, 404 & No & Se repetiría, o es una respuesta válida \\
\bottomrule
\end{tabularx}

Es seguro reintentar porque la consulta a repartidores es \textbf{de solo lectura}. Antes de reintentar
hay una pausa aleatoria de 50 a 150 ms (\emph{jitter}) para que muchas peticiones no vuelvan a la vez.

\subsection{La huella de la petición}

Antes, si alguien reutilizaba una \cod{Idempotency-Key} \textbf{con un pedido distinto}, el sistema
devolvía el pedido viejo sin avisar y el cliente creía que su pedido nuevo existía. Ahora cada pedido
guarda una huella SHA-256 de lo que se pidió (usuario, restaurante y platillos con cantidades, en forma
canónica: ordenados y con repetidos sumados). Con la misma clave: misma huella $\rightarrow$ 200 con el
pedido original; otra huella $\rightarrow$ \textbf{422}. Se comprueba en el atajo inicial y en el caso de
carrera (dos peticiones simultáneas con la misma clave). La columna la agrega \cod{db/07-huella.sql} y
nunca se muestra al cliente.

\subsection{Recuperarse solo y \cod{/estado}}

Los tres servicios usan \cod{ConnectionPool(check=ConnectionPool.check\_connection)}: antes de entregar
una conexión, el pool comprueba que sigue viva; si PostgreSQL se reinició, la descarta y abre otra.
\cod{GET /estado} reúne lo que importa para operar: el estado del circuito y cuántos eventos del
\cod{outbox} faltan por publicar (si sube y no baja, algo pasa con RabbitMQ o con el publicador).

\section{Entorno y cómo repetir las pruebas}

\begin{comando}
cd Proyecto
docker compose down -v
docker compose up --build -d
source ../documentacion/bitacora-4/pruebas/comun.sh      # variables U, R, P, REP y funciones
source ../documentacion/bitacora-4/pruebas/prueba-2.sh   # cualquiera de prueba-0.sh a prueba-7.sh
\end{comando}

\cod{comun.sh} crea un usuario de la zona 3, toma el primer restaurante, su primer platillo y la
repartidora de la zona 3, y define \cod{nuevo\_pedido}, \cod{pagar}, \cod{estado\_pedido},
\cod{aceptar\_tiempo} (código HTTP y tiempo de la aceptación) y \cod{avisos}. La función \cod{paso}
imprime cada comando antes de ejecutarlo, y por eso las salidas de abajo muestran el comando
(\cod{\$ ...}) seguido de lo que contestó el sistema. Las pruebas se corrieron en Git Bash.

\section{Prueba 0: arranque y \cod{/estado}}
\queria Que el sistema levanta completo y \cod{/estado} empieza con el circuito cerrado y sin pendientes.
\evidencia{prueba-0.txt}
\observado Siete contenedores, \cod{/salud} en ok y \cod{/estado} con \cod{CERRADO}, 0 fallos y 0 eventos
pendientes. \hfill\ok

\section{Prueba 1: clave de idempotencia reutilizada con otro pedido}
\queria Que la misma clave con la misma petición devuelve el pedido original, y con otra petición (3
platillos en lugar de 1) se rechaza.
\evidencia{prueba-1.txt}
\observado \cod{201}, \cod{200} y \cod{422} con el mensaje claro. En la base hay \textbf{un solo pedido}
con esa clave, con su huella guardada. \hfill\ok

\section{Prueba 2: repartidores lento y el circuit breaker}
\queria La prueba principal: con repartidores respondiendo en 5 s, las primeras aceptaciones esperan su
tiempo límite y, tras 3 fallos seguidos, las siguientes fallan al instante; al quitar el retraso, el
circuito se cierra solo.
\evidencia{prueba-2.txt}
\observado Los intentos 1 a 3 tardaron \textbf{3.02 s} (\cod{ReadTimeout}, sin reintento) y del 4 al 6,
\textbf{7 ms}: el circuito se abrió y el núcleo dejó de llamar. \cod{/estado} mostró \cod{ABIERTO} y el
log, \cod{CERRADO -> ABIERTO (3 fallos seguidos)}. Al quitar el retraso, el log muestra
\cod{ABIERTO -> SEMIABIERTO} y \cod{SEMIABIERTO -> CERRADO}, y el mismo pedido terminó \cod{ASIGNADO}.

Diferencia con las instrucciones: se esperaba que el primer intento tras quitar el retraso todavía diera
503. Dio 200 porque recrear los contenedores tardó más que la pausa de 10 s, así que esa llamada ya fue la
de prueba. \hfill\ok

\section{Prueba 3: una copia congelada}
\queria Qué pasa si una de las dos copias se congela (\cod{docker pause}): el contenedor sigue existiendo,
así que el DNS de Docker lo sigue entregando a las llamadas internas del núcleo.
\evidencia{prueba-3.txt}
\observado Una mezcla de \cod{200} rápidos (la copia sana) y \cod{503} de 3 s (la congelada). Tres de
esas llamadas cayeron seguidas en la copia congelada, así que \textbf{el circuito se abrió} y las dos
últimas consultas fallaron en 6 ms. Ver la sección~\ref{sec:hallazgos}. \hfill\ok

\section{Prueba 4: repartidores caído por completo}
\queria Que con las dos copias detenidas el núcleo responde 503 rápido, el circuito se abre y, al volver
repartidores, se cierra solo. Se parte con el circuito \cod{CERRADO}.

\subsection{Antes del arreglo}
\lstinputlisting[style=salida]{pruebas/evidencia/prueba-4-antes-del-arreglo.txt}
Cada intento tardaba \textbf{8 s}, más que el propio límite de 3 s del núcleo. Ver la
sección~\ref{sec:hallazgos}.

\subsection{Después del arreglo}
\evidencia{prueba-4.txt}
\observado Los tres primeros intentos tardaron \textbf{1.01 s} (el nuevo límite para resolver el nombre,
sin reintento) y los siguientes \textbf{7 ms} con el circuito abierto. El mensaje explica que el circuito
está abierto. El pedido siguió \cod{PAGADO} y, al volver repartidores, la llamada de prueba salió bien y el
circuito se cerró. (El \cod{outbox\_pendientes: 1} del final es el evento \cod{pedido.asignado} recién
guardado, a punto de publicarse.) \hfill\ok

\figura[0.85]{img/interfaz-circuito-abierto.png}{La misma situación vista en \cod{interfaz.py}: tres 503 de
1 s, luego uno de 8 ms con el mensaje de circuito abierto y la insignia en rojo.}

\section{Prueba 5: el núcleo muere justo después de un pago}
\queria Que un evento ya confirmado no se pierde aunque el núcleo muera (\cod{docker compose kill}, sin
apagado ordenado) antes de publicarlo. Para asegurar que no alcanzara a publicarlo, RabbitMQ estaba
apagado.
\evidencia{prueba-5.txt}
\observado Antes de matar el núcleo, \cod{outbox\_pendientes} era 1. Al volver, bajó a 0 y el aviso de
pago llegó \textbf{una sola vez}: el evento estaba guardado en la base, en la misma transacción que el pago.
También aparece el aviso de retraso, porque el pedido llevó más de 30 s pagado sin repartidor mientras
se hacía la prueba. \hfill\ok

\section{Prueba 6: PostgreSQL se reinicia}
\queria Que tras reiniciar la base todo se recupera sin reiniciar ningún servicio.
\evidencia{prueba-6.txt}
\observado \cod{/salud} respondió ok; se creó y pagó un pedido (núcleo), llegó su aviso (notificaciones)
y se consultaron repartidores (repartidores). \cod{notificaciones} y \cod{traefik} llevaban 8 minutos
encendidos: ninguno se reinició. Durante el reinicio, el publicador y el consumidor registraron su error
y reintentaron solos. \hfill\ok

\section{Prueba 7: el gateway se cae}
\queria Que sin el gateway nadie entra desde afuera, pero los servicios siguen hablándose.
\evidencia{prueba-7.txt}
\observado \cod{000} desde afuera (nada escucha en el 8080), pero el núcleo alcanzó a repartidores por
el DNS de Docker. Al volver Traefik, todo respondió otra vez. \hfill\ok

\section{Hallazgos de la corrida real}\label{sec:hallazgos}

\subsection{1. El nombre que no se resuelve no tenía tiempo límite (corregido)}

\textbf{Qué pasó.} Con repartidores detenido, cada aceptación tardaba 8 s. Dentro del contenedor del
núcleo, \cod{socket.getaddrinfo("repartidores")} tardaba 4.0 s en responder que el nombre no existe (el
DNS de Docker Desktop en Windows consulta hacia afuera antes de rendirse; en Linux, en la etapa 3, fue
casi instantáneo). El tiempo límite de conexión de \cod{httpx} \textbf{no cubre la resolución del
nombre}, y el reintento la duplicaba: $2 \times 4$ s.

\textbf{Arreglo.} En \cod{app/clientes/repartidores.py}, el nombre se resuelve en un hilo aparte con el
mismo límite de conexión (1 s); la llamada HTTP se hace a la IP obtenida, con el encabezado \cod{Host}
original. Si el nombre no se resuelve, no se reintenta: significa que no hay ninguna copia viva. Cada
llamada sigue resolviendo el nombre, así que el reparto entre las dos copias se mantiene (comprobado en
la regresión: 5 y 5 por el gateway, y las llamadas internas llegan a las dos copias).

\textbf{Resultado.} De 8 s a 1 s por intento, y 7 ms con el circuito abierto.

\subsection{2. Con una copia congelada el circuito sí puede abrirse}

El reporte técnico de la etapa predecía que con una copia congelada el circuito seguiría cerrado,
porque los éxitos intercalados reinician la cuenta. En la corrida real, el DNS de Docker no alterna
estrictamente entre copias y tres llamadas seguidas cayeron en la congelada: el circuito se abrió. Es el
comportamiento correcto del breaker (3 fallos seguidos), pero significa que una sola copia congelada
puede cortar por 10 s las llamadas internas aunque la otra copia esté sana. La causa de fondo es la
limitación ya declarada: las llamadas internas por DNS no usan los \emph{healthchecks} del gateway. La
solución sería llamar a repartidores a través del gateway o con un directorio que conozca la salud.

\section{Matriz de caídas}\label{sec:matriz}

\begin{longtable}{@{}p{3.3cm} p{5.1cm} p{4.1cm} p{2.2cm}@{}}
\toprule
\textsf{\textbf{Pieza que falla}} & \textsf{\textbf{Qué pasa}} & \textsf{\textbf{Qué lo resuelve}} & \textsf{\textbf{Evidencia}} \\
\midrule
\endhead
Una copia de repartidores & El sistema sigue con la otra & Traefik la saca del reparto; el DNS deja de resolverla & E3, prueba 5 \\
Las dos copias & Se crean y pagan pedidos; aceptar da 503 en 1 s y luego al instante; aviso de retraso; cancelación y reembolso & Tiempo límite, breaker, reloj de asignación & E3 prueba 6; E4 prueba 4 \\
Repartidores lento & Tres intentos de 3 s y luego fallo instantáneo & Tiempo límite y breaker & E4, prueba 2 \\
Una copia congelada & Llamadas internas de 3 s; el circuito puede abrirse & Parcial: las llamadas internas no usan los \emph{healthchecks} & E4, prueba 3 \\
Notificaciones & Los pedidos siguen; los avisos esperan en la cola & Cola durable, \emph{ack} tardío & E2, prueba 2 \\
RabbitMQ & Los pedidos siguen; los eventos esperan en el \cod{outbox} & Outbox y publicador que reintenta & E2, prueba 6 \\
RabbitMQ se reinicia & La cola y sus mensajes sobreviven & Cola durable, mensajes persistentes, \cod{hostname} fijo & E2, prueba 7 \\
El núcleo muere tras un pago & El evento se publica al volver, una vez & Outbox en la misma transacción & E4, prueba 5 \\
Muerte entre publicar y marcar & El evento se publica otra vez & Consumidor idempotente por \cod{evento\_id} & Diseño; E2, prueba 3 \\
PostgreSQL se reinicia & Errores durante el reinicio; luego todo se recupera solo & Pool que comprueba las conexiones & E4, prueba 6 \\
El gateway & Nadie entra; los servicios siguen hablándose & La llamada interna no pasa por el gateway & E4, prueba 7 \\
Mensaje inválido & Va a la cola de fallidos & \emph{nack} sin reencolar & E2, prueba 4 \\
Petición repetida & Mismo pedido, sin segundo cobro & \cod{Idempotency-Key} con \cod{UNIQUE} & E1, sección 5 \\
Clave reutilizada con otro pedido & 422 & Huella de la petición & E4, prueba 1 \\
Peticiones simultáneas & Una gana; las demás, mismo resultado o 409 & \cod{UPDATE} condicional & E1, secciones 6--8 \\
Nadie paga o nadie acepta & Cancelación; reembolso si había cobro & Relojes de pedido & E1, sección 9 \\
PostgreSQL cae para siempre & Todo se detiene & No resuelto: punto único de falla declarado & --- \\
\bottomrule
\end{longtable}

\section{Resumen de la etapa 4}

\begin{tabularx}{\linewidth}{@{}l X l@{}}
\toprule
\textsf{\textbf{Prueba}} & \textsf{\textbf{Resultado observado}} & \textsf{\textbf{Estado}} \\
\midrule
0. Arranque & 7 contenedores; circuito \cod{CERRADO}; 0 pendientes & \ok \\
1. Clave reutilizada & 201, 200 y 422; un solo pedido & \ok \\
2. Repartidores lento & 3 × 3.02 s, luego 7 ms; \cod{ABIERTO}; se recupera a \cod{CERRADO} & \ok \\
3. Copia congelada & Mezcla de 200 y 503 de 3 s; el circuito se abrió (hallazgo 2) & \ok \\
4. Repartidores caído & Antes: 8 s; después del arreglo: 1 s y luego 7 ms; se recupera & \ok \\
5. El núcleo muere tras un pago & 1 pendiente $\rightarrow$ 0; aviso de pago una sola vez & \ok \\
6. PostgreSQL se reinicia & Los tres servicios se recuperan sin reiniciarse & \ok \\
7. El gateway se cae & 000 desde afuera; el núcleo alcanza a repartidores & \ok \\
Regresión de las etapas 1 a 3 & Todo igual que en sus bitácoras (ver el documento de pruebas) & \ok \\
\bottomrule
\end{tabularx}

\section{Quién hizo qué}

\begin{tabularx}{\linewidth}{@{}p{4.2cm} X@{}}
\toprule
\textsf{\textbf{Integrante}} & \textsf{\textbf{Aportes en la etapa 4}} \\
\midrule
Yukioayax Canek Gabriel Hernández (@YukiGab) & Circuit breaker, reintentos seguros, huella de la petición, \cod{/estado}, pools que se recuperan solos, integración de la etapa 4 sobre la 3, corrida real de las ocho pruebas y de la regresión, arreglo del tiempo límite de la resolución de nombres, \cod{interfaz.py} (aplicación de escritorio) y esta bitácora. \\
Erick Isaac Echeverria Goicochea (@ErickI07) & Panel web de demostración (\cod{/panel/}) integrado con el núcleo; revisión de la etapa. \\
Los dos & Plan de la etapa, matriz de caídas, revisión cruzada y documentación final. \\
\bottomrule
\end{tabularx}

\end{document}
